\section{Combining the finite regulator conversions}
The three finite shifts in \eqref{eq:comp-simplified-real-remainder}
can be combined before evaluating the real contribution. Define
\begin{equation}
 M(X)=\Delta_{\mathrm{rat}}+\Delta_{\log}+\Delta_D,
 \qquad L=\ln(1+X),\quad J=\ln(1+X/2),\quad X>0.
\end{equation}
The shifts include the finite displacement of the double-pole
boundary in the Bose-restoring term. Combining their independently
integrated endpoint moments gives
\begin{equation}
 \begin{aligned}
 M(X)={}&-\frac{P_0}{3X}
 +\frac{2P_1}{X(1+X)}\operatorname{Li}_2(-X)
 -\frac{8(1+X)(X^2-4X-6)}{X^2}L^3\\
 &+\frac{16(1+X)(X^2-4X-6)}{X^2}L\operatorname{Li}_2(-X)
 -\frac{256(1+X)(X+2)^2}{X^3}J\\
 &+\frac{128(1+X)P_2}{3X^3}
 \left[\operatorname{Li}_2\!\left(-\frac{X^2}{4(1+X)}\right)-2JL\right]
 -\frac{P_3L^2+P_4L}{3X^3(1+X)}.
 \end{aligned}
 \label{eq:comp-matching-short}
\end{equation}
Here the polynomials are
\begin{align}
 P_0&=X^5-X^4-8X^3+196X^2+1104X+864,\\
 P_1&=X^5-6X^4-43X^3-32X^2+32X+32,\\
 P_2&=X^3-5X^2-28X-24,\\
 P_3&=3X^7-18X^6-233X^5-96X^4
                 +3008X^3+8240X^2+8000X+2688,\\
 P_4&=X^9+6X^8-18X^7-24X^6+363X^5
                 +1652X^4+412X^3-3984X^2-4704X-1536.
\end{align}
In particular, the explicit $\ln X$ cancels in this sum.
The reduction uses positive-real logarithm identities and
\begin{equation}
 \operatorname{Li}_2\!\left(\frac{X}{1+X}\right)
 =-\operatorname{Li}_2(-X)-\frac12\ln^2(1+X).
\end{equation}
The last identity follows by differentiating both sides and fixing
the zero value at $X=0$. Both sides are real for $X>0$.
The remaining dilogarithm arguments in
\eqref{eq:comp-matching-short} are negative, so their real branches
are fixed throughout the physical interval. Coefficient collection
then requires only rational identities in $X$. Thus the real remainder
can be written more compactly as
\begin{equation}
 \mathcal R_{\mathrm{rem}}(X)
 =F_{\mathrm{rat}}^z+F_W^z+F_{\log}^z+F_D^z
                  +2F_{UW}^{\eta}+M(X).
\end{equation}
The script \nolinkurl{compton_matching_simplify.py} constructs this
sum from the separate endpoint moments and verifies the rational
coefficients exactly. The overall factor $1/[8X(1+X)]$ is still
applied only in \eqref{eq:comp-simplified-master}.
